\section{The degree-eight arithmetic line: exact coordinates and the missing measurement}

We retain $K=\mathbf Q(i)$, $E:y^2=x^3-x$, and the prime $5$.
The free graded Lie algebra on the chosen CM homology basis $x,y$
is embedded in its associative tensor algebra. Derivation degree $m$
means that $x,y$ are sent to Lie words of length $m+1$.
All computations in this part are exact rational computations, with
the integral lattice subsequently localized at $5$.

\begin{statement}{G.1}{Open}{The arithmetic scalar and Lead 15}
The scalar $\lambda$ in the actual arithmetic line
\[
 \frac{d^3\chi^{(4)}_4}{267840}(A_8+\lambda B_8)
\]
is not evaluated here. Lead 15, which asserts $\lambda\in\mathbf Q$,
also remains open. G.4 independently verifies the two ambient vectors
and their coordinates; G.5 verifies that every $5$-adic scalar is
compatible with the universal elliptic and finite CM constraints.
In particular those constraints admit $\lambda=\sqrt6\in\mathbf Q_5$
and do not imply rationality. This is a counterexample to that proposed
method of deducing rationality, not a counterexample for the actual
arithmetic scalar.

The motivic argument in Lead 15 would need a rational subspace in the
\emph{specified} degree-eight derivation space, together with a
comparison identifying its scalar extension with the actual Galois
line. Merely having Betti and de Rham rational structures does not
give that comparison or identify the two structures in these
coordinates. No such theorem is supplied or proved in this return.
\end{statement}

\begin{statement}{G.2}{Open}{The actual integral arithmetic image}
The values of the second coordinate $\nu$ on an actual element of
the eighth Johnson subgroup of $\Delta$ are not calculated here.
Neither is the nonzero ideal which is the image of the first
coordinate $\alpha$. Consequently the arithmetic lattice and its
index in the saturation of its rational line remain open. G.6 gives
the exact index formula in terms of those two missing data and
distinguishes this finite index from the necessarily infinite index
of a rank-one lattice in the rank-two ambient plane.
\end{statement}

\begin{statement}{G.3}{Open}{The repeated-character arithmetic embeddings through degree fourteen}
No complete arithmetic generator list or integral image lattice in
all repeated-character spaces of degrees eight through fourteen is
obtained. The degree-eight measurement still missing in G.1 is
already part of that requested list. A family of universal
derivations, a metabelian character, or dimensions of the lower
bracket images does not determine those generator vectors. The
calculations below are therefore recorded as ambient certificates,
with no promotion to a description of the Galois subgroup.
\end{statement}

\begin{statement}{G.4}{Proved}{An independent exact certificate for the degree-eight plane}
Let $F(x)=y$, $F(y)=0$, and for even $m\geq2$ put
\[
 \epsilon_m(x)=\operatorname{ad}_x^m(y),\qquad
 \epsilon_m(y)=\sum_{j=0}^{m/2-1}(-1)^j
 [\operatorname{ad}_x^j(y),\operatorname{ad}_x^{m-1-j}(y)].
\]
Set
\[
 A_8=\frac{\operatorname{ad}_F^3\epsilon_8}{120},\quad
 B_8=\frac{[\epsilon_4,\operatorname{ad}_F^2\epsilon_4]}2,
 \quad C_8=A_8+\frac{31}{20}B_8.
\]
These vectors annihilate $[x,y]$, have weight $(4,4)$, and change
sign under $x\leftrightarrow y$. For the integral word-coordinate
functionals
\[
 \alpha(D)=-[x^5y^4]D(x),\qquad
 \nu(D)=[x^4yxy^3]D(x)+4[x^5y^4]D(x),
\]
one has
\[
 \begin{array}{c|rrr}
 &A_8&B_8&C_8\\\hline
 \alpha&1&0&1\\
 \nu&-31/20&1&0
 \end{array}.
\]
Writing $W_8=\mathbf Q_5A_8\oplus\mathbf Q_5B_8$, its saturated
lattice of integral Lie derivations is
\[
 W_8\cap\operatorname{Hom}_{\mathbf Z_5}
       (\mathbf Z_5x\oplus\mathbf Z_5y,L_9(\mathbf Z_5x\oplus\mathbf Z_5y))
       =\mathbf Z_5C_8\oplus\mathbf Z_5B_8.
\]
In the free metabelian Lie algebra, written
$\mathbf Qx\oplus\mathbf Qy\oplus\mathbf Q[u,v][x,y]$ with
$u=\operatorname{ad}_x$, $v=\operatorname{ad}_y$ on the commutator
module, the values of $A_8$ are
$u^4v^3[x,y]$ and $u^3v^4[x,y]$, while both values of $B_8$ vanish.
\end{statement}
\begin{proof}
The executable certificate \texttt{code/fgh\_degree8.py} starts from
these formulas, with its own tensor multiplication, commutator, and
derivation implementation; it imports no attached calculation.
It computes all rational associative coefficients, checks the
special equation and the two symmetries, and decomposes each value
in the standard Lyndon basis by triangular subtraction. At each
subtraction the leading word is checked to be Lyndon and its
standard bracket has leading coefficient $1$. The complete vectors
are recorded in \texttt{code/fgh\_degree8.json}.

For the two selected words its coefficients are
\[
 \begin{array}{c|rrr}
 &A_8(x)&B_8(x)&C_8(x)\\\hline
 {[x^5y^4]}&-1&0&-1\\
 {[x^4yxy^3]}&49/20&1&4
 \end{array},
\]
which gives the displayed matrix. The Lyndon coefficients of $B_8$
are integers; those of $C_8$ have denominators only $1,2,4$.
Thus both are integral at $5$. Conversely, an integral vector
$D\in W_8$ has integral word coefficients, so
$\alpha(D),\nu(D)\in\mathbf Z_5$. The coordinate matrix gives
$D=\alpha(D)C_8+\nu(D)B_8$. This proves saturation, without
identifying any lattice vector as arithmetic.

For an independent check of the metabelian assertion, the program
evaluates each Lyndon bracket in the explicit metabelian Lie
algebra above. The bracket of
$ax+by+p(u,v)[x,y]$ and $cx+dy+q(u,v)[x,y]$ has commutator
coordinate
\[
 ad-bc+(au+bv)q-(cu+dv)p.
\]
This elementary formula supplies the stated result and directly
checks that the second direction is invisible metabelianly.
\end{proof}

\begin{statement}{G.5}{Proved}{Universal equations do not select a rational direction}
For every $\lambda\in\mathbf Q_5$, the line
$\mathbf Q_5(A_8+\lambda B_8)$ exponentiates in the completed
positive geometric subgroup of Enriquez's graded universal elliptic
kernel. It has CM character $(4,4)$ and the prescribed interchange
sign. This remains true for an irrational scalar: the root of
$T^2-6$ congruent to $1$ modulo $5$ belongs to $\mathbf Q_5$ and
is not rational. A line in this plane is defined over $\mathbf Q$
if and only if its normalized scalar $\lambda$ is rational.
\end{statement}
\begin{proof}
Enriquez, Proposition 4.19, equation (48), defines
$\delta_{2n}=\epsilon_{2n+2}$ and places it in
$\mathfrak r_{\mathrm{ell}}^{\mathrm{gr}}$.
The paragraph following equation (49) gives
\[
 \exp(\widehat{\mathfrak b}_3^+)\rtimes\mathrm{SL}_2
       \subset R_{\mathrm{ell}}^{\mathrm{gr}},\qquad
 \mathfrak b_3=\langle\mathfrak{sl}_2,\delta_{2n}\rangle
       .
\]
The inclusion is the one explicitly stated there \cite{fghEnriquez}.
The half-sum in (48) pairs into the formula for $\epsilon_m(y)$
in G.4. The source uses the opposite derivation bracket, which does
not change the generated Lie subalgebra. Both $A_8$ and $B_8$
therefore belong to its positive part. Exponentiation of any scalar
linear combination converges in the degree filtration and gives
the asserted group elements to all orders. G.4 checks the CM weight
and interchange sign exactly. The scalar $\sqrt6$ exists by Hensel's
lemma since $1^2-6\equiv0\pmod5$ and $2$ is a unit, and is
irrational since $6$ is not a rational square.

Finally $A_8,B_8$ are rational and independent. A rational vector
spanning the normalized line has a nonzero $A_8$ coefficient and
therefore a rational ratio of its $B_8$ and $A_8$ coefficients.
That ratio is $\lambda$. The converse is immediate. No arithmetic
realization of the irrational candidate is asserted.
\end{proof}

\begin{statement}{G.6}{Proved}{The exact conditional lattice and index formula}
Let $L=\mathbf Q_5(A_8+\lambda B_8)$ be any line in this plane,
put $\delta=\lambda-31/20$, and set
\[
 m=\max\{0,-v_5(\delta)\},\qquad v_5(0)=+\infty.
\]
Then
\[
 L\cap(\mathbf Z_5C_8\oplus\mathbf Z_5B_8)
     =5^m\mathbf Z_5(C_8+\delta B_8).
\]
If $\mathcal A\ne0$ is the compact image of a continuous
homomorphism from a profinite group to this integral line and
$\alpha(\mathcal A)=5^a\mathbf Z_5$, then $a\geq m$ and
\[
 \mathcal A=5^a\mathbf Z_5(C_8+\delta B_8),\qquad
 [L\cap\Lambda:\mathcal A]=5^{a-m},
 \quad\Lambda=\mathbf Z_5C_8\oplus\mathbf Z_5B_8.
\]
Thus this formula applies to the character-$(4,4)$ image of an
eighth Johnson homomorphism whenever its rational line is the
one specified in the prompt. It leaves $a$ and $\delta$ to be
measured on that actual homomorphism.
\end{statement}
\begin{proof}
The coordinates of $t(C_8+\delta B_8)$ are $(t,t\delta)$.
Both are integral exactly when $v_5(t)\geq0$ and
$v_5(t)+v_5(\delta)\geq0$, giving the first formula.
A compact additive subgroup of $\mathbf Q_5$ is closed and stable
under $\mathbf Z_5$, by taking limits of multiplication by integers.
Its nonzero integral image is consequently an ideal
$5^a\mathbf Z_5$. Since $\alpha$ is injective on $L$, it determines
the whole subgroup and the index. The full rank-two quotient
$\Lambda/\mathcal A$ also has a rank-one part, so its cardinal
index is infinite; the finite displayed index is taken only inside
the saturation of $L$.
\end{proof}

\begin{statement}{G.7}{Proved}{A finite nonmetabelian measurement that would decide the scalar}
Let $P$ be the free pro-$5$ group on lifts $a,b$ of the integral
homology basis. An actual automorphism $\sigma$ in its eighth
Johnson subgroup determines $D_8(\sigma)$ by the leading terms of
$\sigma(a)a^{-1}$ and $\sigma(b)b^{-1}$. For every $r\geq1$,
the two coordinates $\alpha,\nu$ modulo $5^r$ can be read from its
action on the characteristic finite quotient supplied by
\[
 (\mathbf Z/5^r)\langle X,Y\rangle/(X,Y)^{10},
 \qquad a\mapsto1+X,\quad b\mapsto1+Y.
\]
On a vector of the arithmetic line with $\alpha\ne0$,
\[
       \lambda=\frac{31}{20}+\frac{\nu}{\alpha}.
\]
This is an explicit test on a genuine arithmetic action; it does
not provide such an action or its measured values.
\end{statement}
\begin{proof}
The completed group algebra of a free pro-$5$ group is the completed
noncommutative power-series algebra under these substitutions.
Its augmentation ideal and the ideal generated by $5^r$ are
characteristic. The quotient algebra is finite and its augmentation
units form a finite $5$-group, as follows from the finite additive
successive powers of the augmentation ideal. Hence the indicated
group quotient and automorphism action are defined.

For an eighth Johnson element, all positive terms of lengths less
than nine in $\sigma(a)a^{-1}-1$ vanish. The degree-nine term is
the Lie element $D_8(\sigma)(x)$ because leading group commutators
expand to the corresponding Lie commutators. The two integral
linear combinations of word coefficients in G.4 therefore compute
$\alpha,\nu$ modulo $5^r$ without division or a logarithm.
The ratio formula is the coordinate identity of G.6. In particular
metabelian data, which kill $B_8$, cannot replace this second
measurement.
\end{proof}

\begin{statement}{G.8}{Known}{The realization hypothesis in the proposed motivic shortcut}
Malkin's \emph{Motivic fundamental groups of CM elliptic curves
and geometry of Bianchi hyperbolic threefolds}, Theorem 11,
proves a depth-two symmetry for Hodge correlators; Theorem 16
constructs the corresponding Bianchi-complex map in that realization.
The transfer discussed immediately before Lemma 3, on printed
pages 12--13, invokes hypothetical regulator injectivity and the
motivic formalism. It therefore supplies no unconditional theorem
identifying the degree-eight $5$-adic arithmetic line with a rational
line in the fixed derivation coordinates used here
\cite{fghMalkin}.
\end{statement}
The primary text has been opened at those statements. This is a
precise limit on what that proposed source proves, not a proof that
the desired comparison cannot be established by another argument.

